\chapter{PENDAHULUAN}
\label{chap:pendahuluan}
\begin{petunjuk}
Bab Pendahuluan secara umum yang dijadikan landasan kerja dan arah kerja penulis tugas akhir, berfungsi mengantar pembaca untuk membaca laporan tugas akhir secara keseluruhan.
\end{petunjuk}

\section{Latar Belakang}
\begin{petunjuk}
Latar Belakang berisi dasar pemikiran, kebutuhan atau alasan yang menjadi ide dari topik tugas akhir. Tujuan utamanya adalah untuk memberikan informasi secukupnya kepada pembaca agar memahami topik yang akan dibahas. Saat menuliskan bagian ini, posisikan anda sebagai pembaca -- apakah anda tertarik untuk terus membaca?
\end{petunjuk}

\section{Rumusan Masalah}
\begin{petunjuk}
Tuliskan rumusan masalah yang menjadi dasar perlunya tugas akhir dibuat. Rumusan masalah yang baik memiliki struktur sebagai berikut:
\begin{enumerate}
  \item Diawali dengan ringkasan dalam satu atau dua kalimat tentang kondisi/situasi dalam bentuk \textit{problem statement} yang ada sekarang terkait dengan topik utama yang dibahas pada tugas akhir.
  \item Pokok persoalan dari kondisi/situasi yang ada, dapat dilihat dari kelemahan atau kekurangannya. Bagian ini merupakan inti dari rumusan masalah.
  \item Elaborasi lebih lanjut yang menekankan hal-hal penting sebagai pokok kajian untuk menyelesaikan persoalan tugas akhir.
  \item Dapat ditambahkan usulan singkat terkait dengan solusi yang ditawarkan untuk menyelesaikan persoalan tugas akhir.
\end{enumerate}
Penting untuk diperhatikan bahwa persoalan yang dideskripsikan pada subbab ini akan dipertanggungjawabkan di akhir pelaksanaan tugas akhir, apakah terselesaikan atau tidak?
\end{petunjuk}

\section{Tujuan dan Ukuran Keberhasilan Pencapaian}
\begin{petunjuk}
Tuliskan tujuan utama dan/atau tujuan detil yang akan dicapai dalam pelaksanaan tugas akhir. Fokuskan pada hasil akhir yang ingin diperoleh setelah tugas akhir diselesaikan, terkait dengan penyelesaian persoalan pada rumusan masalah. Penting untuk diperhatikan bahwa tujuan yang dideskripsikan pada subbab ini akan dipertanggungjawabkan di akhir pelaksanaan tugas akhir apakah tercapai atau tidak. Tuliskan juga ukuran keberhasilannya.
\end{petunjuk}

\section{Batasan Masalah}
\begin{petunjuk}
Tuliskan batasan-batasan kajian yang diambil dalam pelaksanaan tugas akhir dan bukan merupakan batasan implementasi solusi. Tuliskan batasan yang dicakup tersebut dalam kalimat positif. Batasan ini dapat dihindari (tidak perlu ada) jika topik/judul tugas akhir dibuat cukup spesifik.
\end{petunjuk}

\section{Metodologi}
\begin{petunjuk}
Tuliskan semua tahapan yang akan dilalui selama pelaksanaan tugas akhir. Tahapan ini spesifik untuk menyelesaikan persoalan tugas akhir. Tahapan studi literatur tidak perlu dituliskan karena ini adalah pekerjaan yang harus Anda lakukan selama proses pelaksanaan tugas akhir.
\end{petunjuk}
